\newcommand{\rev}[1]{\source{https://github.com/caihongqiu417/RedflagEngine/commit/#1}{#1}}
\section{图集取错颜色、部件重叠与跨窗拖拽}
下面几页把旧场景的问题按当时的线索重走一遍。日志能告诉我们作者最初看到了什么，提交能告诉我们实际改了哪一行。这批旧问题都没有保留下当时的 RenderDoc 帧。

\subsection{Kenney 模型的一部分采到了黑色}
Kenney 模型用一张图集给不同部位上色，导入后却有一片变黑。旧日志沿图片和纹理坐标两头检查：图片加载时已经上下翻转，Assimp 又启用了 \code{aiProcess\_FlipUVs}。修复保留图片加载时的翻转，去掉导入模型时的第二次处理，让模型重新取到图集中对应的色块。

如今 Model 仍未启用这个 Assimp 标志。原 Kenney 资产本次没有重载，其他格式需要根据图片来源和 UV 约定分别验证。
\note{依据：\path{Task/ModelImport_DevLog.md} 第 2 节；导入提交 \rev{3bcd608}。}

\subsection{子网格叠在原点}
多部件模型加载后，明明有几个子网格，画出来却都叠在原点。ProcessNode 递归找到了孩子节点，却漏了从父节点走到当前节点时应乘的变换。修复把父节点与当前节点的矩阵累积到 SubMesh::LocalTransform，绘制时再乘整个实体在场景里的变换。

父子矩阵这条乘法链在提交和当前源码中都能找到。多层节点的专用资产回归仍然缺失；前文的材质用例检查几何有没有被删，未检查层级位置。
\note{依据：同一日志第 6.1 节；提交 \rev{f433b7e}；Model::ProcessNode、Renderer3D::DrawModel。}

\subsection{加入 3D 后，2D 场景全黑}
旧日志里，3D 加入后，原本能画出的 2D 场景变黑了。先查相机。3D 初始化选中了自己的相机数据，随后 2D 即使更新了自己的矩阵，着色器仍可能读 3D 的那份。再查真正绘制的地方，2D 的 Flush 也没有选回自己的着色器和顶点数组。

修复让更新相机时选回对应的数据，提交精灵时选回 2D 程序与顶点数组，并在初始化时恢复贴图数组。两套相机的内存各自存在，错的是绘制当刻所选的对象。后处理双实例测试复现过相似的状态切换，当年的完整 2D 场景没有在本轮重建。
\note{依据：同一日志第 6.2 节；提交 \rev{0e46596}。}

\subsection{跨窗口拖拽中断}
从资源浏览器把图标拖向另一个窗口，鼠标一离开图标，目标窗口就接不到路径。旧日志记下四次尝试：先依赖鼠标悬停，后来改看图标是否活动。拖拽系统接管之后，这两种状态都可能消失。直接调用 BeginDragDropSource 也仍会中断。

最终在鼠标按下时记住资源路径和类型，按住期间由窗口持续用 SourceExtern 发布数据，松开后再清空。ContentBrowserPanel 至今保留这几个成员。这里能对照当时的尝试和当前实现，尚未重测各种跨窗口布局。
\note{依据：同一日志第 3 节 ContentBrowserPanel 的迭代表；当前 ContentBrowserPanel.cpp。}

\newpage
\section{播放时消失的灯和环境图上的黑块}
\subsection{复制场景遗漏灯光和运行时网格}
编辑场景切到运行，方向光、点光源或程序生成的网格可能缺席。查看提交 \rev{3a5c63a}，Scene::Copy 与 DuplicateEntity 的复制清单原先漏了两类光，RenderRuntime 则漏了遍历程序网格。修复分别补齐复制和绘制。

当前源码仍含这两项处理。原提交没有失败画面和自动测试，无法确定当时先看见灯消失，还是先看见模型消失。

\subsection{停止播放后，变换手柄还拿着旧实体}
在播放中选中一个实体，然后点停止。编辑器在这一帧换回编辑场景，但稍后绘制的变换手柄（Gizmo）手里还拿着运行场景的实体。修复把播放时写入选择、操作手柄的入口限制住，防止运行时副本留在编辑状态里。这个过程来自提交说明和代码，全部 Play/Simulate 组合仍未做交互压力测试。
\note{依据：提交 \rev{4c3140b}；ForgeLayer 的鼠标选择与 Gizmo 分支。}

\subsection{环境光图上的黑块从哪里开始}
环境光照的结果是一片黑时，老指南先把漫反射环境图、不同粗糙度下的环境图和 BRDF 查找表分别显示。哪张先出黑块，就从哪张往回看；若读到 NaN，再用颜色标出来源，最后检查输入 HDR。提交 \rev{846e3de} 写的是 NaN 修复，原来那张 HDR 的出错像素和截帧都没留下。

当前代码把 NaN、无穷大等非有限输入改成零，有限亮度限制到 $[0,50]$。这会削弱特别亮的环境光源。旧文中“Inf 求平均必然变 NaN”的解释并不准确，结果还取决于随后做了哪些运算。
\note{依据：\path{Task/IBL_Implementation/IBL_Implementation_Guide.md} 第 10.3、20、21 节；IBLPrecompute.cpp。}

\subsection{粗糙反射的亮点与随相机移动的天空盒}
指南把这些现象分开检查，而不是统一归咎于光照公式。纹理还没有 mip 层就启用 mip 过滤，会使采样条件不完整。修复先使用普通线性过滤，生成层级后再切换。粗糙反射出现亮点，则检查采样与亮源分布，预过滤使用概率密度选择读取的 mip 层。

天空盒用到的立方体在预计算结束后还要继续绘制，因此它的顶点数组要留到 Shutdown。观察矩阵只取旋转，相机平移时天空仍在远处；使用正交投影则跳过天空盒。旧指南还列出六个面的朝向、环境光开关和 BRDF 的几何项。关于这几项，只留有检查建议。

旧照片没有逐项保留下来。前面提到的常量 HDR 测试是后来补做的，用来核对当前预计算。
\note{依据：同一指南第 11、14、18、20、21 节；提交 \rev{846e3de}、\rev{56cacd6}。}

\newpage
\section{灯光读错、泛光重复滤波与环境遮蔽条纹}
\subsection{同一盏灯的颜色和位置读偏了}
旧管线文档记录了一盏读错颜色与位置的灯。C++ 一连放三个 int，只占 12 字节，而着色器把同样的填充写成了 \code{int pad[3]}。按 std140 规则，这类数组的相邻元素间隔 16 字节。于是两端从不同位置开始读下一盏灯。

当时改用三个独立标量作填充，按灯光头 16 字节、每盏灯 64 字节核对。文档与提交保存了写法和改动。原始数据没有保留下来，现在的灯光参数检查只能独立验证当前布局。
\note{依据：\path{Task/RenderPipeline_Doc.md} 第 6 节及 std140 附录；提交 \rev{9e5f194} 的灯光声明。}

\subsection{顶点输出与片元输入在链接时对不上}
顶点着色器送出一个值，片元着色器接收它，两段程序在链接时需要对得上。老管线日志记录了一次中间代码转换后两侧名称不一致的问题。现在的反编译步骤根据相同的 location 生成相同的名字，再编译、链接。

仓库没有保留下来那次链接错误原文和最小失败 Shader，因此不能补写具体驱动报错。当前工程链接与像素测试另有结果，历史名称问题与本轮 GLSL 330 的 uniform 问题也不是同一错误。
\note{依据：同一文档第 11.2 节；OpenGLShader.cpp 中的 renameByLocation。}

\subsection{泛光已经合好，最终画面又滤了一遍}
泛光是先缩小图像提取亮部，再逐级放大、把亮度合回去。查看提交 \rev{dc1457c}，上采样已经把各层合到第一层，最后画到屏幕时又做了一次九点滤波。修复删掉这个额外步骤，直接采样合好的泛光图。当时没有记录误差数值，单凭这段代码也推不出每处边缘都会变硬。

同一提交为 2D Quad 和 Circle 增加颜色到线性空间的近似转换，使写入 HDR 与最终显示转换遵守同一流程。它还把清屏色改成黑色，但没有独立实验说明这项对故障的贡献。本轮后处理测试验证开关、隔离和输出，不是原图的前后视觉对比。
\note{依据：同一提交的 Composite、2D Shader 与 ForgeLayer 差异。}

\subsection{SSAO 重复条纹与采样方向}
屏幕空间环境遮蔽（SSAO）沿表面法线周围选方向取样。如果整面墙使用同一组方向，误差也会沿相同方向排成条纹；若用来生成方向的固定轴正好与法线平行，叉乘还可能变成零。旧实现改为从一张 $4\times4$ 噪声图取方向，去掉沿法线的分量，在每个像素建立自己的切线方向。

提交 \rev{82989dd} 保存了这次改动。当前仍用噪声旋转，没有额外加一遍模糊。增加采样和保留噪声之后的花费，旧记录里缺少同场景的时间或图像指标。随机选方向也仍可能偶遇与法线几乎平行的向量。
\note{依据：\path{Task/SSAO_Implementation_Doc.md} 第 3、6 节；\source{https://learnopengl.com/Advanced-Lighting/SSAO}{LearnOpenGL：SSAO}。}

\newpage
\section{阴影条纹、离地亮缝与间接光拖影}
\subsection{从条纹和亮缝判断该往哪里查}
旧阴影指南把几种看起来相似的瑕疵画在同一张速查表里。黑条纹和离地的亮缝都与深度比较有关，却不能照同一个方向调参数。原日志没有注明表中每一项是否都在 Redflag 复现过，下表据此保留检查步骤。

\noindent\begin{tabularx}{\textwidth}{L{3.3cm}Y}
\toprule
现象 & 日志建议的检查与取舍 \\
\midrule
受光面出现黑条纹 & 对照灯光深度与当前表面，检查深度精度和 bias；增加偏移也可能让阴影离开物体。 \\
阴影与物体间有亮缝 & 检查偏移是否过大，逐项调整，不能同时扩大多个参数后只看最终画面。 \\
场景边缘没有阴影 & 核对光源覆盖范围。扩大正交范围会降低同一张 ShadowMap 的空间精度。 \\
阴影锯齿或模糊 & 区分纹素分辨率与滤波半径。PCF 比较邻近深度后平均遮挡结果。 \\
PCSS 全部显得很软 & 检查光源尺寸与半影计算；原指南没有该参数的自动校准或前后测量。 \\
已有深度但主绘制无阴影 & 先确认主 Shader 读取的是这张图，再检查投影和范围，不先调亮度掩盖问题。 \\
\bottomrule
\end{tabularx}

指南还记载只使用一张方向光阴影图，以及当时跳过透明材质。前者不能由“增加级联”自动变成多灯光阴影。级联解决的是同一方向光在不同距离的覆盖和精度。透明裁剪是否正确则必须分别检查阴影、深度和主绘制，本轮新增验证另列。
\note{依据：\path{Task/Shadow_Implementation_Guide.md} 第 8、9、19、21 节；提交 \rev{56cacd6}；\source{https://learnopengl.com/Advanced-Lighting/Shadows/Shadow-Mapping}{LearnOpenGL：Shadow Mapping}。}

\subsection{用上一帧的光减轻间接光锯齿}

提交 \rev{70f2dd3} 想减轻间接光的锯齿，把本帧间接光与上一帧的颜色分开保存。现在先用当前深度还原表面的位置，再去上一帧对应的位置取颜色，并用周围像素限制旧颜色的范围。原提交没有保留对照截图或测量值，无法从记录中看出镜头停下后还有多少噪点，也看不出移动时物体边缘的拖影有多长。
这条修复开始用上一帧平滑间接光。原先标题里用过 TAA 一词，当前实现仍缺少物体运动信息和前一帧的表面深度对照。新露出的墙面可能把先前遮住它的物体颜色带进来。

\subsection{有白色贴图的黑材质曾被改白}
旧管线文档曾为一批 FBX 写下补丁：有底色图、底色又很暗，就把底色改成白色。作者明确把材质设为黑色时，这条规则同样会触发。早先根据 glass、mirror 等名字移除网格的处理也有类似问题。

现在改查文件是否真的提供了这些属性，保留取值为零的材质。前面用一张白图和一个黑色材质作对照的 glTF 用例，正是在复查这项旧补丁。
\note{依据：\path{Task/RenderPipeline_Doc.md} 第 10 节的 FBX 修复说明；Model.cpp 历史与本轮差异。}

\newpage
\section{精灵残留、视口拉伸与鼠标误操作}
这组提交早于 3D 管线。它们保留了问题和改动，缺少原始截帧及系统复测，因此不能替当前编辑器完整回归。

\subsection{空批次仍触发绘制}
提交 \rev{e80839c} 检查一帧结束时有没有精灵真的提交。原来空批次也上传顶点并 Flush，底层 DrawIndexed 恰好把零理解为使用整个索引缓冲，旧顶点于是有机会又画一次。现在本批索引数大于零才提交。文末保留了当时的代码。这项历史修复缺少原截图与空帧切换的专用回归。

\subsection{视口缩放使相机比例失调}
提交 \rev{783d448} 增加 SceneCamera::SetViewportSize，并在视口变化时通知不固定比例的相机，重新计算投影。差异确认处理的是投影宽高比，不只是修改渲染目标尺寸。本轮没有测试所有停靠布局、DPI 和固定比例组合。

\subsection{操作界面仍触发相机和粒子}
提交 \rev{b23bfd1} 限制视口焦点和悬停状态，并由 ImGuiLayer 根据界面捕获标记阻断事件。这样界面交互不应继续传到场景操作。当前仍保留视口事件控制，但当时的粒子演示与当前编辑器结构不同，不能直接把旧结论推广到所有输入路径。

\subsection{粒子寿命结束仍有残留}
提交 \rev{6d5afd9} 的标题将问题归因于浮点精度，代码实际做了两类改动：到期重置粒子；绘制前过滤极低寿命比例、透明度和尺寸。当前实现保留这些阈值。没有当年的数值轨迹，不能仅凭标题确认唯一根因，也未测阈值对短寿命效果的影响。

\subsection{相机移动方向与棋盘倾斜}
提交 \rev{add9c89} 调整了旧 Sandbox 四方向按键修改相机位置的符号，修正记录中的移动方向。它只说明当时演示的控制约定，不是后续所有相机控制器的方向证明。

棋盘倾斜是一次尝试过、仍未查明的故障。提交 \rev{8f05b29} 改过纹理过滤、纹理坐标越界时的取样方式和初始视口，标题却明确写着问题还在。后来的记录没有给出能对应到这张棋盘的修复。

\subsection{仅凭 GLFW 提交标题还原不了故障}
提交 \rev{1b5f703} 只写了 GLFW fix，差异主要是子模块版本变化。没有现场和复现条件，无法从这两条信息还原故障。旧日志里还有“排序减少 DrawCall”的说法：当时排序主要让相邻模型少切换状态，独立网格仍逐个绘制。次数示意表没有测量帧率。
\note{历史入口均可用 \code{git show <提交>} 复查。这里列的是已检索到的记录，不保证未写进版本库的故障也能完整恢复。}
